% Ban thu: chi giu cac cau co trong CD2

%@ id=TT26.tholoc-l2.003 | trang=1 | so=2.1 | de=II | vitri=CD2.5,C3.9.6 | muc=9 | tt=0 | thieu=loigiai
\begin{bt}
Cho hai biểu thức $A = \dfrac{\sqrt{x}-2}{2\sqrt{x}+1}$ và $B = \dfrac{\sqrt{x}+1}{\sqrt{x}+2} + \dfrac{\sqrt{x}}{\sqrt{x}-2} + \dfrac{x-\sqrt{x}-2}{4-x}$ với $x \ge 0; x \ne 4$.
1. Tính giá trị của biểu thức $A$ khi $x = 4 + 2\sqrt{3}$.
2. Rút gọn biểu thức $B$.
3. Tìm giá trị nhỏ nhất của biểu thức $P = x - 3B.(\sqrt{x}-2) + 5$.
\loigiai{
}
\end{bt}
